\chapter{Formulación del problema de sintonización como optimización}
\label{ch:formulacion-optimizacion}

\section{Introducción del capítulo}
Cuando la sintonización PID deja de verse como una regla fija y se entiende como una búsqueda
de desempeño óptimo, el problema cambia de naturaleza. Ya no se trata solo de aplicar una
tabla, sino de encontrar el conjunto de ganancias que mejor responde a una planta concreta
bajo perturbaciones específicas. Esta perspectiva es particularmente útil en sistemas aislados,
donde el margen de error operativo es pequeño y las consecuencias de una mala regulación se
manifiestan de forma inmediata.

\section{Objetivo del capítulo}
El objetivo de este capítulo es definir formalmente la sintonización del PID como un problema
de optimización restringida. Para ello se identifican las variables de decisión, se formula una
función objetivo consistente con la dinámica de frecuencia y se delimitan las restricciones
físicas y numéricas que condicionan la búsqueda.

\section{Variables de decisión}
Las variables de decisión son las ganancias del controlador: $K_p$, $K_i$ y $K_d$. Cada una
modifica un aspecto distinto de la respuesta, por lo que el espacio de búsqueda es
tridimensional y presenta interacciones no triviales entre rapidez, error acumulado y
amortiguamiento. En términos prácticos, cualquier algoritmo de optimización debe proponer
tripletas de ganancias candidatas, simular su efecto sobre la planta y conservar aquellas que
ofrecen mejor desempeño.

\section{Función objetivo}
La función objetivo adoptada en esta obra se basa en el criterio ITAE, definido como
\[
J = \int_0^{t_{\text{sim}}} t\,|\Delta f(t)|\,dt,
\]
donde $\Delta f(t)$ representa la desviación de frecuencia. Este índice resulta conveniente
porque penaliza con fuerza los errores que persisten en el tiempo, favoreciendo respuestas que
no solo reaccionan rápido, sino que se estabilizan sin oscilaciones prolongadas. Así, la
optimización se orienta hacia un comportamiento útil desde la operación real del sistema.

\section{Restricciones}
La búsqueda de ganancias no puede realizarse sobre un dominio ilimitado. Deben imponerse
rangos admisibles para $K_p$, $K_i$ y $K_d$, así como condiciones de estabilidad, límites del
actuador y criterios de viabilidad numérica. En aplicaciones hidráulicas también es necesario
considerar saturación del gobernador, variaciones de caudal y el riesgo de exigir respuestas
demasiado agresivas que, aunque matemáticamente atractivas, serían inconvenientes para la
vida útil del sistema.

\section{Espacio de búsqueda}
El espacio de búsqueda puede interpretarse como el conjunto de todas las tripletas
$(K_p,K_i,K_d)$ que satisfacen las restricciones previamente definidas. Dentro de ese dominio,
algunas regiones producen respuestas lentas pero estables; otras generan oscilaciones o
inestabilidad; y unas pocas concentran soluciones prometedoras. Justamente ahí radica la
utilidad de los métodos metaheurísticos: recorrer de manera eficiente un espacio irregular sin
necesidad de gradientes exactos ni modelos convexos.

\section{Modelo de planta y entorno de simulación}
Como planta de referencia se considera una microred híbrida hidro-solar basada en una turbina
de flujo cruzado Michell--Banki de 200 kW. El modelo incorpora la dinámica del gobernador y
la inercia de la columna de agua, responsable del comportamiento de fase no mínima asociado
al golpe de ariete. La validación se plantea en MATLAB/Simulink mediante escenarios de
incremento de carga y caída de irradiancia, lo que permite evaluar cada candidato PID en
condiciones comparables y técnicamente exigentes.

\section{Cierre del capítulo}
Una vez formulado el problema, la cuestión central deja de ser si existe una buena sintonía y
pasa a ser cómo encontrarla de manera eficaz. El capítulo siguiente introduce la lógica general
de los métodos metaheurísticos, que constituyen el marco de búsqueda empleado en esta obra.
